\textbf{1. Kiến thức cần nhớ}
\begin{itemize}
\item \textbf{HÌNH CẦU}: Tạo thành khi quay nửa hình tròn tâm $O$ bán kính $R$ một vòng quanh đường kính cố định.
\item \textbf{Diện tích mặt cầu} bán kính $R$ (đường kính $d = 2R$):
$$\boxed{S = 4\pi R^2 = \pi d^2}$$
\item \textbf{Thể tích hình cầu}:
$$\boxed{V = \dfrac{4}{3}\pi R^3 = \dfrac{1}{6}\pi d^3}$$
\end{itemize}

\textbf{2. Phương pháp giải các dạng}

%% DANG 9H0?32-1
\textbf{Dạng 1. Diện tích mặt cầu, thể tích hình cầu}
\begin{enumerate}
\item Xác định bán kính $R$ hoặc đường kính $d$.
\item Áp dụng trực tiếp công thức diện tích $S = 4\pi R^2$ và thể tích $V = \dfrac{4}{3}\pi R^3$.
\item Khi biết diện tích $S$, tìm bán kính: $\boxed{R = \sqrt{\dfrac{S}{4\pi}}}$.
\item Khi biết thể tích $V$, tìm bán kính: $\boxed{R = \sqrt[3]{\dfrac{3V}{4\pi}}}$.
\end{enumerate}

%% DANG 9H0?32-2
\textbf{Dạng 2. Bài toán thực tế về hình cầu (Bài 5 đề Tuyển sinh 10)}
\begin{enumerate}
\item Các vật thể thực tế: Quả bóng đá, quả địa cầu, bồn chứa khí ga hình cầu, quả cam hình cầu, viên bi thép\dots
\item Độ dày vỏ quả cầu: Ruột cầu có bán kính $r = R - d_{\text{vỏ}}$.
\item Tính khối lượng: $m = D \cdot V$ (với $D$ là khối lượng riêng).
\item Bài toán phối hợp: Thả viên bi hình cầu vào bình nước hình trụ $\rightarrow$ thể tích nước dâng lên bằng thể tích viên bi:
$$\Delta V = S_{\text{đáy trụ}} \cdot h_{\text{dâng}} = V_{\text{cầu}} \implies \pi R_{\text{trụ}}^2 \cdot h_{\text{dâng}} = \dfrac{4}{3}\pi R_{\text{cầu}}^3.$$
\end{enumerate}
